\documentclass[10pt,a4paper]{amsart}
\usepackage[margin=19mm]{geometry}
\usepackage{amsmath,amssymb,mathtools}
\usepackage[scr=rsfs,cal=euler]{mathalfa}
\usepackage{xfrac}
\usepackage[unicode,hidelinks,bookmarks=false]{hyperref}
\newcommand{\ie}{i.e.\ }
\newcommand{\cf}{cf.\ }
\newcommand{\resp}{respectively\ }
\newcommand{\Subsection}{\subsection}
\providecommand{\nobreakdash}{\nobreak\hbox{-}}
\setlength{\emergencystretch}{2em}
\multlinegap0pt
\usepackage{tikz,xcolor}
\usetikzlibrary{calc,through,backgrounds,shapes,matrix,math,arrows,trees}
\usepackage[nameinlink]{cleveref}
\newtheorem{theorem}{Theorem}[section]
\theoremstyle{definition}
\crefname{defi}{Definition}{Definitions}
\crefname{prop}{Proposition}{Propositions}
\crefname{coro}{Corollary}{Corollaries}
\usepackage{aliascnt}
\newaliascnt{lemma}{theorem}
\newtheorem{lemma}[lemma]{Lemma}
\aliascntresetthe{lemma}
\crefname{lemma}{Lemma}{Lemmas}
\newaliascnt{prop}{theorem}
\newtheorem{prop}[prop]{Proposition}
\aliascntresetthe{prop}
\crefname{prop}{Proposition}{Propositions}
\newaliascnt{coro}{theorem}
\newtheorem{coro}[coro]{Corollary}
\aliascntresetthe{coro}
\crefname{coro}{Corollary}{Corollarys}
\newaliascnt{defi}{theorem}
\newtheorem{defi}[defi]{Definition}
\aliascntresetthe{defi}
\crefname{defi}{Definition}{Definitions}
\newaliascnt{rema}{theorem}
\newtheorem{rema}[rema]{Remark}
\aliascntresetthe{rema}
\crefname{rema}{Remark}{Remarks}
\definecolor{darkblue}{rgb}{0.0,0,0.7} 
\newcommand{\darkblue}{\color{darkblue}} 

\newcommand{\defn}[1]{\emph{\darkblue #1}}

\newcommand{\pop}{\mathsf{Pop}}
\newcommand{\wo}{w_\circ}
\newcommand{\U}{{W'}}
\newcommand{\uu}{{w'}}

\DeclareMathOperator{\spn}{span}
\DeclareMathOperator{\NC}{NC}
\DeclareMathOperator{\Des}{Des}
\DeclareMathOperator{\mov}{Mov}
\DeclareMathOperator{\fold}{fold}
\DeclareMathOperator{\im}{im}
\DeclareMathOperator{\unfold}{unfold}
\DeclareMathOperator{\cyc}{cyc}
\DeclareMathOperator{\aexc}{aexc}
\DeclareMathOperator{\Aexc}{Aexc}
\DeclareMathOperator{\Sym}{Sym}
\DeclareMathOperator{\Weak}{Weak}


\DeclarePairedDelimiter\abs{\lvert}{\rvert}

\newcommand*{\mk}{\mkern -1mu}
\newcommand*{\Mk}{\mkern -2mu}
\newcommand*{\mK}{\mkern 1mu}
\newcommand*{\MK}{\mkern 2mu}

\hypersetup{urlcolor=purple, linkcolor=blue, citecolor=red}


\newcommand*{\romanenumi}{\renewcommand*{\theenumi}{\roman{enumi}}}
\newcommand*{\Romanenumi}{\renewcommand*{\theenumi}{\Roman{enumi}}}
\newcommand*{\alphenumi}{\renewcommand*{\theenumi}{\alph{enumi}}}
\newcommand*{\Alphenumi}{\renewcommand*{\theenumi}{\Alph{enumi}}}

\title[Coxeter inverse-element pop-tsack orbit]{Coxeter Pop-Tsack Torsing: original Theorem1.3}
\author{Colin Defant and Nathan Williams}
\date{Source: Algebraic Combinatorics5(3)(2022),559--581; DOI10.5802/alco.226}
\begin{document}\maketitle
\noindent\textit{Editorial scope.} The counted unit is original Theorem1.3; all other retained results and the editable folding diagram are uncounted core. The source\textquotesingle s complete alternative inverse-power proof is retained. The earlier alternative orbit-length route, redundant restatements, type-specific classifications and illustrative permutahedron example are omitted. No original mathematical expression or assumption is repaired.
\section{Original operator and counted theorem}
\setcounter{equation}{2}
Let $T$ be the set of all reflections of $W$, and fix a Coxeter element $c$ of $W$ (a regular element of order $h$).  Let $\leq_T$ denote the absolute order on $W$. There is a general ``dual'' philosophy for finite Coxeter groups~\cite{bessis2003dual} in which the set $S$ of simple reflections is replaced by the set $T$ of all reflections, the longest element $\wo$ is replaced by the Coxeter element $c$, and the (left) weak order is replaced by the noncrossing partition lattice $\NC(W,c)$ (the interval $[e,c]$ in absolute order).

Define the \defn{noncrossing projection}
\begin{align}\label{EqNCProj}
\pi_T(\cdot,c)&:W \to \NC(W,c) \nonumber\\
\pi_T(w,c) &= \bigvee_{t \leq_T w}^{\NC(W,c)} t.
\end{align}

\setcounter{theorem}{1}
\begin{defi}
\label{def:dual_pop_stack}
The \defn{Coxeter pop-tsack torsing operator}\footnote{We have interchanged the roles of $S$ and $T$.} is the map
\begin{gather*}
 \pop_T(\cdot,c):W\to W \nonumber \\
\pop_T(w,c) = w \cdot \pi_T(w,c)^{-1}.
\end{gather*}
\end{defi}

When $c$ is understood, we may abbreviate $\pi_T(w,c)$ and $\pop_T(w,c)$ as $\pi_T(w)$ and $\pop_T(w)$, respectively; we establish in~\Cref{sec:general_properties} many structural properties of $\pop_T$, including that its orbit structure does not depend on $c$.  
By construction, the elements $w\in W$ such that $\pop_T(w)=e$ are exactly the noncrossing partitions---hence, the number of such elements is the \defn{$W$-Catalan number} 
\[
\left|\NC(W,c)\right|=\prod_{i=1}^n \frac{h+d_i}{d_i},
\]
 where $d_1,\ldots,d_n$ are the degrees of $W$.

\begin{theorem}
\label{thm:cinverseorbit}
The forward orbit of $c^{-1}$ under $\pop_T(\cdot,c)$ is 
\[
O_{\pop_T}(c^{-1})=\{c^{-1}, c^{-2},\ldots, c^{-(h-1)}, e\}.
\]
 The element $c^{-1}$ has no preimages under $\pop_T$. Moreover, if $2\leq i\leq h-1$, then the only preimage of $c^{-i}$ under $\pop_T$ is $c^{-(i-1)}$. 
\end{theorem}

\section{Preliminaries}
Throughout this article, we assume that $W$ is a finite irreducible Coxeter group with a set $S$ of simple reflections. The \defn{rank} of $W$ is $|S|$. Given $g,w\in W$, we write $w^g$ for the conjugate $gwg^{-1}$ of $w$ by $g$. Let $T=\{s^w:s\in S, w\in W\}$ denote the set of \defn{reflections} of $W$. A \defn{$T$-word} for an element $w\in W$ is a word $t_1\cdots t_r$ over the alphabet $T$ that is equal to $w$ when considered as an element of $W$. The \defn{reflection length} of $w$, denoted $\ell_T(w)$, is the minimum length of a $T$-word for $w$. The \defn{absolute order} on $W$, denoted $\leq_T$, is defined by saying $v\leq_T w$ if $\ell_T(wv^{-1})=\ell_T(w)-\ell_T(v)$. Since $T$ is closed under conjugation by elements of $W$, we have $v\leq_T w$ if and only if $\ell_T(v^{-1}w)=\ell_T(w)-\ell_T(v)$. 

It is often useful to consider the \defn{geometric representation} of $W$ on a vector space $V\cong\mathbb R^{|S|}$. We refer the reader to \cite{armstrong2009generalized, BjornerBrenti} for more thorough discussions of this representation; here, we simply mention the basic results that we will need. We denote the action of an element $w\in W$ on a vector $\alpha\in V$ by $w\cdot\alpha$. For every $t\in T$, there is a hyperplane $H_t\subseteq V$ such that $t$ acts via the linear transformation that reflects through $H_t$. Furthermore, $w\cdot H_t=H_{t^w}$ for every $w\in W$ and $t\in T$. Define the \defn{moved space} of $w$ to be $\mov(w)=\im(w-1)$, the image of the linear transformation $w-1:V\to V$ ($1$ is the identity transformation). Brady and Watt \cite{BradyWatt} showed that if $t\in T$ and $w\in W$, then \begin{equation}\label{EqBradyWatt}
\mov(t)\subseteq\mov(w)\text{ if and only if }t\leq _T w.
\end{equation}

The ring $\Sym(V_\mathbb{C}^*)^W$ of $W$-invariant polynomials is a polynomial ring $\mathbb{C}[f_1,\ldots,f_n]$ over a set of invariant polynomials $f_i$ that may be chosen homogeneous and whose degrees $d_1 \leq \cdots \leq d_n$ are uniquely determined.  We write $h=d_n$ for the \defn{Coxeter number} of $W$.

\looseness-1
A \defn{standard Coxeter element} of $W$ is an element obtained by multiplying the simple reflections together in some order.  A \defn{regular element} of $W$ is an element with an eigenvector that lies in the complement of the reflecting hyperplanes.  Following~\cite{reiner2017non}, we define a \defn{Coxeter element} of $W$ to be a regular element of order $h$.  For Weyl groups, Coxeter elements coincide with conjugates of standard Coxeter elements.  
More generally, a \defn{reflection automorphism} of $W$ is a group automorphism that preserves the set of reflections of $W$.  By~\cite[Proposition~1.4]{reiner2017non}, the reflection automorphisms of $W$ act transitively on its Coxeter elements---thus, all Coxeter elements have the same order $h$.

We assume throughout this article that we have fixed a Coxeter element $c$ of $W$. A \defn{noncrossing partition} is an element $v\in W$ such that $v\leq_T c$. The interval $[e,c]$ in the absolute order is called the \defn{noncrossing partition lattice} and denoted $\NC(W,c)$; it is indeed a lattice. This is the only type
of lattice we will consider in the rest of the article, so all meets (denoted by $\bigwedge$) and joins (denoted by $\bigvee$) will be with respect to this lattice. Furthermore, the definitions of the noncrossing projection $\pi_T$ and the Coxeter pop-tsack torsing operator $\pop_T$ from \eqref{EqNCProj} and \Cref{def:dual_pop_stack} are made with respect to this 
fixed Coxeter element. Since the reflection automorphisms of $W$ act transitively on the Coxeter elements, different Coxeter elements give rise to isomorphic noncrossing partition lattices.

\section{Dihedral Groups}
\label{sec:dihedral}

As the dihedral case will be useful when analyzing the general case, we immediately dispense with it in this short section.  The dihedral group $I_2(h)$ has $2h$ elements: the identity $e$, $h$ reflections, a Coxeter element $c$, and then the elements $c^{-1},c^{-2},\ldots,c^{-(h-1)}$.  The identity, reflections, and Coxeter element are all noncrossing partitions, and they all map to the identity under a single application of $\pop_T$.  For each $1\leq i\leq h-1$, the element $c^{-i}$ requires exactly $h-i$ iterations of $\pop_T$ to reach the identity.

\setcounter{figure}{4}
\section{General Properties of \texorpdfstring{$\pop_T$}{Pop}}
\label{sec:general_properties}

\subsection{Absolute monotonicity}

Given $x\in W$, let $W_x$ be the subgroup of $W$ generated by the reflections lying weakly below $x$ in the absolute order. The subgroups $W_x$ for $x\in \NC(W,c)$ are called the \defn{noncrossing parabolic subgroups} of $W$. As discussed in \cite{ReadingNoncrossing}, the noncrossing parabolic subgroups of $W$ form a lattice under the containment order, and the map $x\mapsto W_x$ gives an isomorphism from $\NC(W,c)$ to this lattice. 

\begin{lemma}
\label{lem:monotone}
For every $x \in W$, we have $\pi_T(\pop_T(x))\leq_T \pi_T(x)$.
\end{lemma}
\begin{proof}
The parabolic subgroup generated by the reflections lying weakly below $\pi_T(x)$ in absolute order is $W_{\pi_T(x)}$. All of the reflections lying weakly below $x$ also lie weakly below $\pi_T(x)$, so $x$ is an element of $W_{\pi_T(x)}$. Of course, $\pi_T(x)$ also belongs to $W_{\pi_T(x)}$. This means that $\pop_T(x)=x\pi_T(x)^{-1}$ is in $W_{\pi_T(x)}$, so all of the reflections weakly below $\pop_T(x)$ are in $W_{\pi_T(x)}$. It follows that 
$W_{\pi_T(\pop_T(x))}$ is contained in $W_{\pi_T(x)}$, and this proves the result. 
\end{proof}

The following proposition gives a condition for elements to have an ``isolated'' forward orbit under $\pop_T$. It states that the only possible preimage under $\pop_T$ of an element $w \in W$ with $\pi_T(w)=c$ is $wc$.

\begin{prop}
\label{prop:preimages}
If $x,w\in W$ are such that $\pop_T(x)=w$ and $\pi_T(w)=c$, then $x=wc$.
\end{prop}

\begin{proof}
Suppose $x$ is such that $\pop_T(x)=w$.  By \Cref{lem:monotone}, $c\leq_T\pi_T(x)$, so we must have $\pi_T(x)=c$. However, this implies that $x=\pop_T(x)\pi_T(x)=wc$.  
\end{proof}

\subsection{Equivariance under reflection automorphisms}

In this section, we make a point of decorating $\pi_T$ and $\pop_T$ by the Coxeter element $c$ in order to differentiate different choices of Coxeter elements.  We show a sort of equivariance of $\pop_T$ with respect to reflection automorphisms, which establishes that the orbit structure of $\pop_T$ does not depend on the choice of Coxeter element $c$.  Thus, the Coxeter pop-tsack torsing operators defined with respect to different Coxeter elements are dynamically equivalent to each other, so none of our results about $\pop_T$ depend on the choice of $c$. This means that we have the freedom to choose any Coxeter element we want when we work with specific Coxeter groups, and it will often be helpful to use a specific choice that makes the resulting combinatorics easiest to describe.

\begin{prop}  Let $\phi$ be a reflection automorphism of $W$.
For $w \in W$, we have
\begin{itemize}
\item $\phi(\pi_T(w,\phi^{-1}(c))) = \pi_T(\phi(w),c)$, and
\item $\phi(\pop_T(w,\phi^{-1}(c))) = \pop_T(\phi(w),c)$.
\end{itemize}
\label{prop:pop_props}
\end{prop}

\begin{proof}
For the first point, observe that the map $w \mapsto \phi(w)$ is a lattice isomorphism from $\NC(W,c)$ to $\NC(W,\phi(c))$.  Furthermore, for each reflection $t$, we have $t \leq_T w$ if and only if $\phi(t) \leq_T \phi(w)$.  Therefore, 
\[
\phi\left(\pi_T(w,\phi^{-1}(c))\right) = \phi\left(\bigvee_{t \leq_T w}^{\NC(W,\phi^{-1}(c))} t\right) = \bigvee_{t \leq_T w}^{\NC(W,c)} \phi(t) = \bigvee_{t \leq_T \phi(w)}^{\NC(W,c)} t = \pi_T(\phi(w),c).\]  


For the second point, we use the first point to check that 
\begin{align*}\pop_T(\phi(w),c) &= \phi(w) \pi_T(\phi(w),c)^{-1} = \phi(w) \left(\phi\left(\pi_T(w,\phi^{-1}(c))\right)\right)^{-1}\\ & = \phi\left(w\pi_T(w,\phi^{-1}( c))^{-1}\right) = \phi\left(\pop_T(w,\phi^{-1}(c))\right) .\qedhere
\end{align*}
\end{proof}

As conjugation by $c$ is a special kind of reflection automorphism that restricts to a lattice isomorphism of $\NC(W,c)$, we can specialize the preceding result to conclude the following.

\begin{coro} For $w \in W$, we have
\begin{itemize}
    \item  $\pi_T(w,c)^c = \pi_T(w^c,c)$, and
    \item $\pop_T(w,c)^c = \pop_T(w^c,c)$. 
\end{itemize}
\end{coro}


Since reflection automorphisms act transitively on the set of Coxeter elements, we further obtain:
\begin{coro}
\label{cor:orbit_structure}
The orbit structure of $\pop_T$ does not depend on the choice of Coxeter element.
\end{coro}

\subsection{Folding and the Coxeter plane}
\label{sec:folding}

 


\begin{defi}
Let $(W,S)$ and $(\U,S')$ be Coxeter systems such that $W$ and $W'$ are finite and irreducible. Suppose there is a surjective map $\fold:S\to S'$ such that for each $s'\in S'$, the elements of $\fold^{-1}(s')$ all commute with each other. If there is an injective homomorphism $\unfold:\U\to W$ with $\unfold(s')=\prod_{s\in\fold^{-1}(s')}s$ for all $s'\in S'$, then we say that $W$ \defn{folds} to $W'$ and that $W'$ \defn{unfolds} into $W$. 
\end{defi}

We remark that our definition of folding is more general than the standard definition that uses an automorphism of the Dynkin diagram. 

\Cref{fig:folding} illustrates several examples of folding. If $W$ folds to $\U$ and $c'$ is a Coxeter element of $\U$, then $\unfold(c')$ is a Coxeter element of $W$. Hence, the Coxeter number is preserved by folding.


A useful example of folding goes through the Coxeter plane.  Any finite irreducible Coxeter group $W$ with Coxeter number $h$ folds to the dihedral group $I_2(h)$ of order $2h$, as we now explain.  Since the Coxeter--Dynkin diagram of $W$ is a tree (hence, a bipartite graph), the set $S$ of simple reflections of $W$ naturally separates into two sets $S_+$ and $S_-$, where the simple reflections within each set all commute with each other.  Writing $c_\pm$ for the product of the reflections in $S_\pm$ so that $c = c_+ c_-$ is a Coxeter element of $W$, $\langle c_+,c_-\rangle$ is the dihedral group $I_2(h)$.  Each of $c_+$ and $c_-$ acts as a reflection on the \defn{Coxeter plane}, a plane spanned by the real and imaginary parts of an eigenvector for $c$ with eigenvalue $e^{2\pi i/h}$.


\begin{figure}[htbp]
    \centering
    \bgroup
\def\arraystretch{2}
\begin{tabular}{|ccc|c|}
\hline
     $W$ & $\U$ & $h$ & Diagrammatic folding  \\ \hline
     $A_{2n-1}$ & $B_n$ & $2n$\vrule height20pt depth20pt width0pt & \raisebox{-0.5\height}{\begin{tikzpicture}[scale=.5]
       \filldraw[black] (0,.5) circle (2pt) -- (1,0) circle (2pt) -- (2,0) circle (2pt);
        \filldraw[black,dotted] (2,0) circle (2pt) -- (3,0) circle (2pt);
        \filldraw[black] (3,0) circle (2pt) -- (4,0) circle (2pt);
        \filldraw[black] (0,.5) circle (2pt) -- (1,1) circle (2pt) -- (2,1) circle (2pt);
        \filldraw[black,dotted] (2,1) circle (2pt) -- (3,1) circle (2pt);
        \filldraw[black] (3,1) circle (2pt) -- (4,1) circle (2pt);
        \filldraw[black,double] (0,-.5) -- (1,-.5);
        \filldraw[black] (0,-.5) circle (2pt);
        \filldraw[black] (1,-.5) circle (2pt);
        \filldraw[black] (1,-.5) circle (2pt) -- (2,-.5) circle (2pt);
        \filldraw[black,dotted] (2,-.5) circle (2pt) -- (3,-.5) circle (2pt);
        \filldraw[black] (3,-.5) circle (2pt) -- (4,-.5) circle (2pt);
      \end{tikzpicture}}\\ 
     $D_n$ & $B_{n-1}$ & $2n-2$\vrule height0pt depth20pt width0pt & \raisebox{-0.5\height}{\begin{tikzpicture}[scale=.5]
     \filldraw[black] (0,1) circle (2pt) -- (1,.5) circle (2pt);
       \filldraw[black] (0,0) circle (2pt) -- (1,.5) circle (2pt) -- (2,.5) circle (2pt);
        \filldraw[black,dotted] (2,.5) circle (2pt) -- (3,.5) circle (2pt);
        \filldraw[black] (3,.5) circle (2pt) -- (4,.5) circle (2pt);
        \filldraw[black,double] (0,-.5) -- (1,-.5);
        \filldraw[black] (0,-.5) circle (2pt);
        \filldraw[black] (1,-.5) circle (2pt);
        \filldraw[black] (1,-.5) circle (2pt) -- (2,-.5) circle (2pt);
        \filldraw[black,dotted] (2,-.5) circle (2pt) -- (3,-.5) circle (2pt);
        \filldraw[black] (3,-.5) circle (2pt) -- (4,-.5) circle (2pt);
      \end{tikzpicture}} \\ 
     $E_6$ & $F_4$ & $12$\vrule height0pt depth25pt width0pt & \raisebox{-0.5\height}{\begin{tikzpicture}[scale=.5]
       \filldraw[black] (-1,.5) circle(2pt) -- (0,.5) circle (2pt) -- (1,0) circle (2pt) -- (2,0) circle (2pt);
        \filldraw[black] (0,.5) circle (2pt) -- (1,1) circle (2pt) -- (2,1) circle (2pt);
        \filldraw[black,double] (0,-.5) -- (1,-.5);
        \filldraw[black] (-1,-.5) circle (2pt) -- (0,-.5) circle (2pt);
        \filldraw[black] (0,-.5) circle (2pt);
        \filldraw[black] (1,-.5) circle (2pt);
        \filldraw[black] (1,-.5) circle (2pt) -- (2,-.5) circle (2pt);
      \end{tikzpicture}}\\ 
     $E_8$ & $H_4$ & 30\vrule height0pt depth20pt width0pt & \raisebox{-0.5\height}{\begin{tikzpicture}[scale=.5]
       \filldraw[black] (0,1.5) circle(2pt) -- (1,1) circle (2pt) -- (2,1) circle (2pt) -- (3,1) circle (2pt);
        \filldraw[black] (0,.5) circle (2pt) -- (1,0) circle (2pt) -- (2,0) circle (2pt) -- (3,0) circle (2pt);
        \filldraw[black] (0,.5) -- (1,1);
        \filldraw[black] (0,-.5) circle (2pt) -- (1,-.5) circle (2pt) -- (2,-.5) circle (2pt) -- (3,-.5) circle (2pt);
        \node at (.5,-1) {\scriptsize 5};
      \end{tikzpicture}}\\
      \hline
\end{tabular}
\egroup
    \caption{Examples of folding of finite Coxeter groups.}
    \label{fig:folding}
\end{figure}

In this subsection, we write $T$ and $T'$ for the sets of reflections in $W$ and $W'$, respectively. 

\begin{lemma}
Let $\U$ be a finite irreducible Coxeter group that unfolds into $W$. Let $c'$ be a bipartite Coxeter element of $\U$, and let $c=\unfold(c')$ be the corresponding Coxeter element of $W$. Then $\unfold(\NC(\U,c'))$ is a sublattice of $\NC(W,c)$.
\label{lem:lattice_embedding}
\end{lemma}
\begin{proof}
Suppose that $w'\in W'$ and $w' \leq_{T'} c'$ so that $\uu$ can be found as a prefix of $c'$. Then $\unfold(\uu)$ can be found as a prefix of $c$, so $\unfold(\uu) \leq_T c$.  By the same reasoning (replacing now $c'$ by any element of $\NC(\U,c')$), we have that $\unfold(\NC(\U,c'))$ is a subposet of $\NC(W,c)$.

We now check that $\unfold(\NC(\U,c'))$ is a sublattice of $\NC(W,c)$ (that is, that $\unfold(\NC(\U,c'))$ is closed under the meet and join operations of $\NC(W,c)$).  We know by~\cite{brady2008non} that a noncrossing partition $w$ is characterized by the set of reflections that lie below it: $w = \bigvee_{t \leq_T w} t$. Since 
\[
w \wedge v = \bigvee\{t : t \leq_T w \text{ and } t \leq_T v\},
\]
 it follows that $\unfold(\NC(\U))$ is a meet-sublattice. Indeed, for fixed $t'\in T'$ and $r \in \fold^{-1}(t')$, we have $r \leq \unfold(w') \wedge \unfold(v')$ if and only if $r \leq_T \unfold(w')$ and $r \leq_T \unfold(v')$. This occurs if and only if $t' \leq_{T'} w'$ and $t' \leq_{T'} v'$, which occurs if and only if $t' \leq_{T'} w' \wedge v'$, which occurs if and only if $r \leq_T \unfold(w' \wedge v')$. 


The Kreweras complement maps $K:\NC(W,c)\to \NC(W,c)$ and $K':\NC(W',c')\to \NC(W',c')$, defined by $K(w)= cw^{-1}$ and $K'(w')=c'(w')^{-1}$, are lattice anti-isomorphisms \cite{armstrong2009generalized}. We have $K\circ\unfold=\unfold\circ K'$ because $\unfold$ is a homomorphism that sends $c'$ to $c$. This shows that $K$ preserves the set $\unfold(\NC(\U,c'))$, so $\unfold(\NC(\U,c'))$ is also a join-sublattice of $\NC(W,c)$.
\end{proof}

\setcounter{subsection}{3}
\subsection{The inverse Coxeter element orbit}
\label{sec:inverse_c_orbit}

\setcounter{theorem}{11}
\begin{lemma}\label{Lem:joinofcorbit}
If $t\in T$ is a reflection, then $c=\bigvee\limits_{k=0}^{h-1} t^{c^{k}}$.
\end{lemma}

\begin{proof}
Since $W$ folds to $I_2(h)$ via the Coxeter plane, we can consider the Coxeter element $\unfold(c')$ of $W$, where $c'$ is a Coxeter element of $I_2(h)$. Each Coxeter element of $W$ can be obtained by applying a reflection automorphism to $\unfold(c')$, so it suffices to prove the lemma when $c=\unfold(c')$.

Conjugation by $c$ induces an automorphism of $\NC(W,c)$; therefore, if a set $O \subseteq \NC(W,c)$ is invariant under conjugation by $c$, then the element $x=\bigvee_{w \in O} w$ is invariant under conjugation by $c$:
	\[
x^c = \left(\bigvee_{w \in O} w\right)^c = \bigvee_{w \in O} w^c = \bigvee_{w \in O} w = x.
\]
It is a classical result that the centralizer of $c$ is the cyclic group generated by $c$ (see \cite[Proposition~30]{Carter}), so we must have $x=c^i$ for some integer $i$.  Therefore, we are left to prove that the only powers of $c$ in $\NC(W,c)$ are $c$ and $e$. This will follow if we can show that the only powers of $c'$ in $\NC(I_2(h),c')$ are $c'$ and $e$.  But this is clear since the absolute length of any power of $c'$ (a rotation) is $0$ or $2$ (see also~\cite[Lemma~12.2]{bessis2015finite}).
\end{proof}

\begin{prop}
For each integer $1\leq i\leq h-1$, we have $\pop_T(c^{-i})=c^{-(i+1)}$. Thus, the forward orbit of $c^{-1}$ under $\pop_T$ is $\{c^{-1},c^{-2},\ldots,c^{-(h-1)},e\}$, which has size~$h$.
\label{prop:forward_cinv}
\end{prop}

\begin{proof}
Fix $1\leq i\leq h-1$. We claim that if $t\in T$ is such that $t \leq_T c^{-i}$, then $t^c \leq_T c^{-i}$.  According to \eqref{EqBradyWatt}, we have $\mov(t) \subseteq \mov(c^{-i})$.  Take $\alpha_t$ to be a vector orthogonal to the hyperplane $H_t$ so that $\mov(t) = \spn\{\alpha_t\}$.  Then $\alpha_t\in\mov(c^{-i})$, so there exists a vector $\beta\in V$ such that $(c^{-i}-1)\cdot\beta = \alpha_t$. Applying $c$ to each side of this equation shows that 
\[
(c^{-i}-1)\cdot(c\cdot\beta)=c\cdot((c^{-i}-1)\cdot\beta)=c\cdot\alpha_t.
\]
 Thus, $\spn\{c\cdot\alpha_t\}\subseteq\mov(c^{-i})$. Since $c\cdot H_t=H_{t^c}$, we know that $c\cdot\alpha_t$ is orthogonal to $H_{t^c}$. Thus, $\mov(t^c)=\spn\{c\cdot\alpha_t\}$. Appealing to \eqref{EqBradyWatt} again, we find that $t^c\leq_T c^{-i}$. 

Note that there must exist some $t\in T$ with $t \leq_T c^{-i}$ since $c^{-i} \neq e$. It follows from the preceding paragraph that $t^{c^k}\leq_T c^{-i}$ for all $0\leq k\leq h-1$. Therefore, we can use \Cref{Lem:joinofcorbit} to see that 
\[
c=\bigvee\limits_{k=0}^{h-1} t^{c^{k}}\leq_T\bigvee_{t'\leq_T c^{-i}}t'=\pi_T(c^{-i}).
\]
 Consequently, $\pi_T(c^{-i})=c$, so $\pop_T(c^{-i})=c^{-i}c^{-1}=c^{-(i+1)}$. 
\end{proof}

\noindent\textit{Original closing preimage argument.}
\begin{proof}
\looseness-1
  By \Cref{prop:preimages} and \Cref{prop:forward_cinv}, the only possible preimage of $c^{-i}$ under $\pop_T$ is $c^{-(i-1)}$.  Furthermore, $c^{-(i-1)}$ is actually a preimage of $c^{-i}$ unless $i=1$.
\end{proof}

This completes the proof of \Cref{thm:cinverseorbit}.

\begin{thebibliography}{99}
\bibitem[3]{armstrong2009generalized} Drew Armstrong, \emph{Generalized noncrossing partitions and combinatorics of Coxeter groups}, Mem. Amer. Math. Soc.202, no.949, American Mathematical Society, Providence, RI,2009.
\bibitem[7]{bessis2003dual} David Bessis, ``The dual braid monoid'', \emph{Ann. Sci. \'Ecole Norm. Sup.}(4)\textbf{36}(2003),no.5,647--683.
\bibitem[8]{bessis2015finite} David Bessis, ``Finite complex reflection arrangements are $K(\pi,1)$'', \emph{Ann. of Math.}(2)\textbf{181}(2015),no.3,809--904.
\bibitem[9]{BjornerBrenti} Anders Bj\"orner and Francesco Brenti, \emph{Combinatorics of Coxeter groups}, Graduate Texts in Mathematics231, Springer, New York,2005.
\bibitem[11]{BradyWatt} Thomas Brady and Colum Watt, ``$K(\pi,1)$'s for Artin groups of finite type'', \emph{Geom. Dedicata}\textbf{94}(2002),225--250.
\bibitem[12]{brady2008non} Thomas Brady and Colum Watt, ``Non-crossing partition lattices in finite real reflection groups'', \emph{Trans. Amer. Math. Soc.}\textbf{360}(2008),no.4,1983--2005.
\bibitem[16]{Carter} Roger W. Carter, ``Conjugacy classes in the Weyl group'', \emph{Compositio Math.}\textbf{25}(1972),1--59.
\bibitem[23]{ReadingNoncrossing} Nathan Reading, ``Noncrossing partitions, clusters and the Coxeter plane'', \emph{S\'em. Lothar. Combin.}\textbf{63}(2010),PaperB63b(32pages).
\bibitem[24]{reiner2017non} Victor Reiner, Vivien Ripoll and Christian Stump, ``On non-conjugate Coxeter elements in well-generated reflection groups'', \emph{Math. Z.}\textbf{285}(2017),no.3-4,1041--1062.
\end{thebibliography}
\end{document}
